\documentclass[12pt, letterpaper]{article}

%-------Packages---------
\usepackage{amssymb,amsfonts}
\usepackage{mathptmx}
\usepackage{amsthm}
\usepackage{graphicx}
\usepackage[all,arc]{xy}
\usepackage{enumerate}
\usepackage{bbm}
\usepackage{dsfont}
\usepackage{mathrsfs}
\usepackage{tikz-cd}
\usepackage{setspace}
\usepackage{import}
\usepackage[utf8]{inputenc}
\usepackage{hyperref}
\usepackage{tikz}
\usepackage{float}
\usepackage[dvipsnames]{xcolor}
\usepackage[nocheck]{fancyhdr}
\usepackage[nointegrals]{wasysym}

\usepackage[left=1in,top=1in,right=1in,bottom=1in]{geometry}

\usepackage{mathtools}
\usepackage{setspace}
\usepackage{cleveref}
\usepackage{multicol}
\usepackage{mathpazo}
\usepackage{tcolorbox}
\tcbuselibrary{theorems,skins,breakable}
\usepackage{xparse}


\usepackage{xcolor, ifthen, tcolorbox}
\tcbuselibrary{theorems,skins,breakable}
% Theorem System
% The following environments are provided:
%   New Problem:        \begin{problem} ...
%   New Question Header:   \qh (then followed by subproblems)
%   Subproblem:     \begin{subproblem} ...
%   Problem Proof:  \begin{problemproof} ...
%   Proof:          \\begin{pf}{color} ...
%   Remark:         \rmk (sentence), \rmkb (block)

% Define custom colors
\definecolor{thmscol}{HTML}{F4565B} %For Problems
\definecolor{lemscol}{HTML}{FE844C} %For Lemmes
\definecolor{prpscol}{HTML}{00A7EF} %For proposition
\definecolor{corscol}{HTML}{23DAFF} %For corrolaries
\definecolor{rmkscol}{HTML}{6DEEC5} %For remarks
\definecolor{defscol}{HTML}{18C991} %For definitions
\definecolor{exmscol}{HTML}{7DB800} %For examples
\definecolor{clmscol}{HTML}{165c58} %For claims

%Change between dark(black)/light(white) mode
\newcommand{\colormode}{white}
\ifthenelse{\equal{\colormode}{white}}
      {\newcommand{\TextColor}{black}}
      {\newcommand{\TextColor}{white}}
\newcommand{\pbcolormain}{thmscol!80!\colormode}
\newcommand{\pbcolorsecond}{thmscol!38!\colormode}


% ============================
% Problem Counter
% ============================
\newcounter{subproblemcounter}
\newcounter{problemcounter}

\newcommand{\q}{
    \setcounter{subproblemcounter}{0}%
    \stepcounter{problemcounter} % Increment problem counter
    \color{\TextColor}Problem \arabic{problemcounter}: % Display the problem counter
}

\newcommand{\stepsub}{
    \stepcounter{subproblemcounter}%
    \color{\TextColor}(\alph{subproblemcounter})
}
% ============================
% Problem
% ============================
\newtcolorbox{pb}[1]{
  enhanced,
  frame hidden,
  titlerule=0mm,
  toptitle=1mm,
  bottomtitle=1mm,
  fonttitle=\bfseries\large,
  coltitle=black,
  colbacktitle=\pbcolormain,
  colback=\pbcolorsecond,
  title={\q #1},
  coltext=\TextColor
}

\newtcolorbox{spb}[1]{
  enhanced,
  frame hidden,
  titlerule=0mm,
  toptitle=1mm,
  bottomtitle=1mm,
  fonttitle=\bfseries\large,
  coltitle=black,
  colbacktitle=\pbcolormain,
  colback=\pbcolorsecond,
  title={\stepsub #1},
  coltext=\TextColor,
  attach boxed title to top left={xshift=3mm,yshift=-2mm},
  boxed title style={
    colback=\pbcolormain,
    colframe=\pbcolormain,
  }
}

\newtcolorbox{pbh}[1]{
    enhanced,
    frame hidden,
    titlerule=0mm,
    toptitle=1mm,
    bottomtitle=1mm,
    fonttitle=\bfseries\large,
    coltitle=black,
    colbacktitle=\pbcolormain,
    colback=\pbcolorsecond,
    title={\q #1},
    boxrule=0pt,
    interior hidden,
    attach boxed title to top left={xshift=3mm,yshift=-2mm},
    boxed title style={
        colback=\pbcolormain,
        colframe=\pbcolormain,
    }
}

\newenvironment{problem}[1]{%
  \newpage
  \begin{pb}{#1}
    }{
  \end{pb}
}

\newenvironment{subproblem}[1]{%
  \begin{spb}{#1}
    }{
  \end{spb}
}

\newenvironment{problemproof}{%
    \begin{pf}{\pbcolormain}
        }{
    \end{pf}
}

\newcommand{\qh}[1][]{
    \newpage
    \begin{pbh}{#1}\end{pbh} % Display the problem counter
}

% ============================
% Claim
% ============================

\newtcbtheorem[number within=problemcounter]{claim}{Claim}
{
    enhanced,
    frame hidden,
    titlerule=0mm,
    toptitle=1mm,
    bottomtitle=1mm,
    fonttitle=\bfseries\large,
    coltitle=black,
    colbacktitle=clmscol!50!\colormode,
    colback=clmscol!20!\colormode,
}{clm}

\NewDocumentCommand{\clm}{m+m}{
    \begin{claim*}{#1}{}
        #2
    \end{claim*}
}

\newenvironment{clmpf}{
	{\noindent{\it \textbf{Proof.}}
	\setlength{\parindent}{0pt}
	}
	\tcolorbox[blanker,breakable,left=5mm,parbox=false,
    before upper={\parindent15pt},
    after skip=10pt,
	borderline west={1mm}{0pt}{clmscol!50!\colormode}]
}{
    \textcolor{clmscol!50!\colormode}{\hbox{}\nobreak\hfill$\blacksquare$}
    \endtcolorbox
}

\NewDocumentCommand{\clmp}{m+m+m}{
    \begin{claim*}{#1}{}
        #2
    \end{claim*}

    \begin{clmpf}
        #3
    \end{clmpf}
}


% ============================
% Proof
% ============================

\newenvironment{pf}[1]{%
  \def\pfcolor{#1}% Store the color in a macro
  \noindent{\it \textbf{Proof.}}%
  \tcolorbox[blanker,breakable,left=5mm,parbox=false,%
    before upper={\parindent15pt},%
    after skip=10pt,%
    borderline west={1mm}{0pt}{\pfcolor},%
    coltext=\TextColor
  ]%
  \setlength{\parindent}{0pt}%
}{%
  \textcolor{\pfcolor}{\hbox{}\nobreak\hfill$\blacksquare$}%
  \endtcolorbox%
}

\newtcolorbox{solution}{
  blanker,
  breakable,
  left=5mm,
  parbox=false,%
  before upper={\parindent15pt},%
  after skip=10pt,%
  borderline west={1mm}{0pt}{\pbcolormain},%
  coltext=\TextColor
}


% ============================
% Remark
% ============================


\NewDocumentCommand{\rmk}{+m}{
    {\it \color{rmkscol!100!white}#1}
}

\newenvironment{remark}{
    \par
    \vspace{5pt}
    \begin{minipage}{\textwidth}
        {\par\noindent{\textbf{Remark.}}}
        \tcolorbox[blanker,breakable,left=5mm,
        before skip=10pt,after skip=10pt,
        borderline west={1mm}{0pt}{rmkscol!80!\colormode},
        coltext=\TextColor
        ]
}{
        \endtcolorbox
    \end{minipage}
    \vspace{5pt}
}

\NewDocumentCommand{\rmkb}{+m}{
    \begin{remark}
        #1
    \end{remark}
}

\usepackage[all]{xypic}

% Define a new command for problems
\renewcommand{\headrule}{\color{\TextColor}\hrulefill}
\newcommand{\C}{\mathbb{C}}
\newcommand{\R}{\mathbb{R}}
\newcommand{\Q}{\mathbb{Q}}
\newcommand{\Z}{\mathbb{Z}}
\newcommand{\N}{\mathbb{N}}
\renewcommand{\P}{\mathbb{P}}
\newcommand{\F}{\mathbb{F}}
\newcommand{\contr}{\Rightarrow \Leftarrow}
\newcommand{\s}{\(\,\)}
\renewcommand{\t}[1]{\text{#1}}
\newcommand{\ang}[1]{\left\langle #1 \right\rangle}
\newcommand{\coker}{\mathrm{coker}}

% Meta data
\setstretch{1.2}

\begin{document}
\color{\TextColor}
\pagecolor{\colormode}
\setlength{\parindent}{0pt}
\pagestyle{fancy}
\fancyhf{}
\fancyhead[LOE]{\color{\TextColor}\slshape{\textbf{Vincent Ferrara}}}
\fancyhead[ROE]{\color{\TextColor}\thepage}

\begin{problem}{}
    Let $\Omega$ be a field and let $K$ and $L$ be subfields of $\
Omega$.
\begin{equation*}
\xymatrix{
& \Omega\ar@{-}[dl]\ar@{-}[dr] &\\
K\ar@{-} & & L\ar@{-}\\
}
\end{equation*}
Let $KL$ denote the smallest subfield of $\Omega$ containing both $K$ and $L$. $KL$
is called the \textit{compositum} of $K$ and $L$. Give an explicit description of
the elements in $KL$.
\end{problem}
\begin{solution}
\setlength{\parindent}{0pt}
Let $S=\{\sum_{i=1}^n k_il_i \mid n \in \N,\; a_i \in K,\; b_i \in L\}$

Let $s,t \in S$ then
\[s=\sum_{i=1}^n a_ib_i \quad \t{and} \quad t=\sum_{j=1}^m c_jd_j\]

Thus,
\[st=\sum \sum (a_ic_j)(b_id_j) \in S\]

Also, $st=0$ only if $s$ or $t$ is zero since $K$,$L$ are subfields of $\Omega$. Thus, $S$ is also multiplicative.

Thus, $S$ is an intergral domain. Thus, consider the fraction field of $S$. Call this $F$.

First, $K \subset F$ and $L \subset F$ since $1 \cdot k \in F$ and $1 \cdot l \in F$ for all $k \in K$ and $l \in L$.

Thus, $KL \subset F$.

Let $E$ be a subfield of $\Omega$ s.t. $K \subset E$ and $L \subset E$.

Since $E$ is a field it is closed under addition and multiplication. Thus, for any $a_i \in K$ and $b_i \in L$, each product $a_ib_i \in E$. Thus, er get that any finite sum of such products is also in $E$. Thus, $S \subset E$.

Since $E$ is a filed, $E$ is also closed under inverses except for 0. Thus, if $t \in S$ and $t\neq 0$ then $t^{-1} \in E$.

Thus, $st^{-1} \in E$ for all $s,t \in S$ and $t \neq 0$.

Therefore, $F \subset E$.

Thus, $F = KL$.
\end{solution}

\begin{problem}{}
    A class of extensions is said to be distinguished if it satisfies the following
three properties:
\begin{enumerate}
\item[i.] Given a tower of extensions $L/K/F$:
\begin{equation*}
\xymatrix{
L\ar@{-}[d]\\
K\ar@{-}[d]\\
F
}
\end{equation*}
the extension $L/F$ is in the class if and only if $L/K$ and $K/F$ are in the
class.
\item[ii.] Given a diagram of field extensions as below,
\begin{equation*}
\xymatrix{
& \Omega\ar@{-}[d]\\
& KL\ar@{-}[dl]\ar@{-}[dr] &\\
K\ar@{-}[dr] & & L\ar@{-}[dl]\\
& F &
}
\end{equation*}
if $K/F$ is in the class, then $KL/L$ is in the class.
\item[iii.] If $K/F$ and $L/F$ are in the class, and $K$ and $L$ are contained in a
field $\Omega$, then the compositum $KL/F$ is in the class.
\end{enumerate}
Note that (iii) follows formally from (i) and (ii).
\end{problem}

\begin{subproblem}{}
    Show that finite extensions form a distinguished class.
\end{subproblem}
\begin{problemproof}
    Property $(i)$ follows directly from what we have proved in class.

    Assume $K/F$ is finite.

    Thus, $K=F(a_1,\dots,a_n)$ for $a_1,\dots,a_n \in K$

    Thus, consider $L(K)$. This is the smallest field of $L$ adjoined with elements of $K$. Thus, $KL=L(K)=L(F(a_1,\dots,a_n))$.

    Since $F \subset L$ adding $F$ does not do anything to $L$. Thus, $L(F(a_1,\dots,a_n))=L(a_1,\dots,a_n)$. Thus, $KL/L$ is finite.
\end{problemproof}

\begin{subproblem}{}
    Show that algebraic extensions form a distinguished class.
\end{subproblem}
\begin{problemproof}
    Property $(i)$ follows from class.

    Assume $K/F$ is algebraic.

    Consider $x \in KL$. Thus, there is a finite number of elements of $K$ that make up $x$. Call these $(a_1,\dots,a_n)$. Consider $L(a_1,\dots,a_n)/L$. Since $K/F$ is algebraic this means that $a_1,\dots,a_n$ satisfy polynomials with coefficients in $F \subset L$. Thus, these also work for $L$.

    Therefore, $L(a_1,\dots,a_n)/L$ is algebraic. Since $x \in L(a_1,\dots,a_n)$. This implies that $x$ is algebraic over $L$. Thus, $KL/L$ is algebraic.
\end{problemproof}

\qh
\begin{subproblem}{}
    Show directly that $X^3-2$ is irreducible over the field $\Q(\omega)
$, by using the arithmetic of the ring $\Z[\omega]$.
\end{subproblem}
\begin{problemproof}
    Assume for contradiction that there exists a nontrivial factorization
    \[X^3-2=(X-\alpha)(X^2+BX+C)\]
    s.t. $\alpha,B,C \in \Z[\omega]$

    Then, $\alpha^3=2$.

    Thus, $N(\alpha)^3=N(2)=4$.

    This implies that $N(\alpha)=4^{1/3} \notin \Z$

    Thus, this is a contradiction so, $X^3-2$ is irreducible in $\Z[\omega][X]$. Therefore, by Gauss' Lemma since $\Z[\omega]$ is a UFD and $\Q[\omega]$ is its fraction field. $X^3-2$ is also irreducible over $\Q[\omega]$.
\end{problemproof}

\begin{subproblem}{}
    Find the irreducible polynomial (over $\Q$) satisfied by $\alpha^2 +
\omega \alpha $.
\end{subproblem}
\begin{problemproof}
    Consider $\Q(\alpha^2+\alpha\omega)$

    Consider $(\alpha(\alpha+\omega))^3=6+6\omega(\alpha^2+\alpha\omega)$

    Since $\alpha^2+\alpha\omega \in \Q(\alpha^2+\alpha\omega)$ and $6 \in \alpha^2+\alpha\omega$. This implies that $\omega \in \Q(\alpha^2+\alpha\omega)$.

    Let $y=\alpha^2+\alpha\omega$

    Thus, $ay=2+\omega\alpha^2=2+\omega(y-\alpha\omega)=2+\omega y-\omega^2\alpha \implies \alpha(y+\omega^2)=2+\omega y \implies \alpha=(2+\omega y)/(y+\omega^2)$

    Thus, $\alpha \in \Q(y)$

    Thus, $\Q(y)=\Q(\alpha,\omega)$

    $f(x)=x^6+6x^4-12x^3+36x^2-36x+36$

    $f(y)=0$ you can check

    Assume $f(x)$ is reducible over $\Q$. Then this means that $y$ will be in one of those factors that has degree less than 6. This is a contradiction since we showed that $\Q(y)/\Q$ has degree 6.

    Thus, $f(x)$ is irreducible over $\Q$.
\end{problemproof}

\qh
\begin{subproblem}{}
    Show that the polynomial $X^4 - X - 1$ is irreducible over $\Q$
\end{subproblem}
\begin{problemproof}
    Assume $X^4-X-1$ is not irreducible in $\Q$. Since $X^4-X-1$ is primitive by Gauss's lemma it is also reducible over $\Z$.
    
    Consider $X^4-X-1=(X-r)(X^3+aX^2+bX+c)$ for $r,a,b,c \in \Z$.

    Thus, $X^4-X-1=X^4+(a-r)X^3+(b-ar)x^2+(c-br)x-cr$

    $X^3: a-r=0 \implies a=r$

    $X^2: b-ar=0 \implies b=r^2$

    $X: c-br=-1 \implies c=r^3-1$

    $cr=1 \implies r^4-r=1$

    $r(r^3-1)=1$

    Thus, $r$ divides both sides. This implies that $r=\pm 1$.

    Both of these solutions do not work.

    Thus, consider $X^4-X-1=(X^2+aX+B)(X^2+cX+d)=X^4+(a+c)X^3+(ac+b+d)X^2+(ad+bc)X+bd$

    $X^3: a=-c$

    $X^2: b+d=a^2$

    $X: a(d-b)=-1$

    $bd=-1$

    This implies that $b,d$ are units thus $b=1,d=-1$ or $b=-1,d=1$.

    In either case we get

    $b+d=0=a^2 \implies a=0$.

    This is a contradiction since this would imply that $0(d-b)=0=-1$

    Thus, $X^4-X-1$ is irreducible over $\Q$.
\end{problemproof}

\begin{subproblem}{}
    Let $\alpha$ be a root of the equation $X^4 - X - 1$ in $\C$. Find
the irreducible polynomial satisfied by $\alpha^2$ over $\Q$.
\end{subproblem}
\begin{problemproof}
    $f(X)=X^4-2X^2-X+1$

    $f(\alpha^2)=\alpha^8-2\alpha^4-\alpha^2+1=(\alpha^4-1)^2-\alpha^2$

    Since $\alpha^4-\alpha-1=0 \implies \alpha=\alpha^4-1$

    $(\alpha^4-1)^2-\alpha^2=\alpha^2-\alpha^2=0$.

    Since this is primitive we can check if it is irreducible over $\Z$.

    Assume it is not irreducible over $\Z$, then

    If $X^4-2X^2-X+1=(X-r)(X^3+aX^2+bX+c)=X^4+(a-r)X^3+(b-ar)X^2+(c-br)X-cr$

    $X^3: a=r$

    $X^2: b=-2+r^2$

    $X: c=r^3-2r-1$

    $r^4-2r^2-r=-1 \implies r(r^3-2r-1)=-1$ Thus $r$ divides $-1$. Thus, $r=-1,1$.

    Then $(1)^4-2(1)^2-1+1=-1 \neq 0$ and $(-1)^4-2(-1)^2+1+1=1 \neq 0$

    If $X^4-2X^2-X+1=(X^2+aX^+b)(X^2+cX^2+d)=X^4+(a+c)X^3+(ac+b+d)X^2+(ad+bc)X+bd$

    $X^3: a=-c$

    $X^2: b+d=a^2-2$

    $X: a(d-b)=-1$

    $bd=1$

    Thus, $b=d=-1,1$

    This is a contradiction since $a(d-b)=-1 \implies 0=-1$

    Thus, $f(X)$ is irreducible over $\Z$ thus also over $\Q$ since it is primitive.
\end{problemproof}

\begin{problem}{}
    Let $K = \Q (\alpha, i )$ where $\alpha$ is a root of the polynomial $X^4 -
2$.
\end{problem}

\begin{subproblem}{}
    Find the degree of $K/\Q$.
\end{subproblem}
\begin{problemproof}
    First this polynomial is irreducible by Eisenstein's Criterion using $p=2$.

    Thus, $\Q(\alpha)/ \Q$ has degree 4.

    Let $\alpha=\sqrt[4]{2}$. Then we know that $\Q(\alpha)$ is a real number field, and so $X^2-1$ which only has complex roots is also not irreducible.

    Thus, $\Q(\alpha,i)/\Q(\alpha)$ has degree 2.

    Therefore, $K/ \Q$ has degree 8.
\end{problemproof}

\begin{subproblem}{}
    Find the irreducible polynomial over $\Q$ satisfied by $\alpha + i$.
\end{subproblem}
\begin{problemproof}
    First consider $\Q(\alpha+i)$. Let $z=\alpha+i$

    $(z-i)^4=2 \implies z^4-4z^3i-6z^2+4zi+1=2 \implies i=(z^4-6z^2-1)/(4(z^3-z))$

    Thus, $i \in \Q(z)$. Thus, $\alpha \in \Q(z)$. Thus, $\Q(z)=\Q(\alpha,i)$

    $f(x)=x^8+4x^6+2x^4+28x^2+1$

    $f(\alpha+i)=4+16i\alpha^3-48i\alpha^2+64i\alpha+24-16i\alpha^3+38\alpha^2+64i\alpha-28=0$

    Assume that $f(x)$ is reducible over $\Q$. This is a contradiction since this would imply that $\alpha+i$ would be in a proper factor that has degree less than 8. Howeverm $\Q(\alpha+i)/\Q$ has degree $8$.

    Thus, $f(x)$ is irreducible.
\end{problemproof}
\end{document}